\chapter*{Abstract}
\addcontentsline{toc}{chapter}{Abstract}

The Cubilox Mk1 is a self-balancing cube that can stand on one of its edges or on one of its
corners (vertices). Three reaction wheels, each driven by a brushless Nidec~24H motor, are mounted
inside a 3D-printed frame. By accelerating and decelerating these wheels, the cube produces
reaction torques that keep its centre of mass above the edge or the corner it is standing on.
An ESP32 microcontroller estimates the orientation of the cube from an inertial measurement unit
(IMU) and runs two controllers: a single-axis controller for the edge and a three-dimensional
controller for the vertex, which separates the tilt motion from the rotation about the vertical
axis.

The project was designed and built entirely by the author: the mechanical parts were designed in
Fusion~360 and printed in PLA, the electronics were soldered on a prototype board, and the firmware
was written from scratch. It was inspired by the Cubli of ETH Zurich~\cite{gajamohan2012cubli,
gajamohan2013cubli} and, more directly, by the self-balancing cube of remrc~\cite{remrc_cube}.

This document describes the working principle, the mechanical and electrical design, the firmware
and its control laws, how to build, flash, calibrate and tune the cube, the results that were
achieved, and the lessons learned during development. The final version balances on its edge and,
after a calibration by hand, on its vertex for several minutes.
